\IfFileExists{stacks-project.cls}{%
\documentclass{stacks-project}
}{%
\documentclass{amsart}
}

% For dealing with references we use the comment environment
\usepackage{verbatim}
\newenvironment{reference}{\comment}{\endcomment}
%\newenvironment{reference}{}{}
\newenvironment{slogan}{\comment}{\endcomment}
\newenvironment{history}{\comment}{\endcomment}

% For commutative diagrams we use Xy-pic
\usepackage[all]{xy}

% We use 2cell for 2-commutative diagrams.
\xyoption{2cell}
\UseAllTwocells

% We use multicol for the list of chapters between chapters
\usepackage{multicol}

% This is generally recommended for better output
\usepackage{lmodern}
\usepackage[T1]{fontenc}

% For cross-file-references
\usepackage{xr-hyper}

% Package for hypertext links:
\usepackage{hyperref}

% For any local file, say "hello.tex" you want to link to please
% use \externaldocument[hello-]{hello}
\externaldocument[introduction-]{introduction}
\externaldocument[conventions-]{conventions}
\externaldocument[sets-]{sets}
\externaldocument[categories-]{categories}
\externaldocument[topology-]{topology}
\externaldocument[sheaves-]{sheaves}
\externaldocument[sites-]{sites}
\externaldocument[stacks-]{stacks}
\externaldocument[fields-]{fields}
\externaldocument[algebra-]{algebra}
\externaldocument[brauer-]{brauer}
\externaldocument[homology-]{homology}
\externaldocument[derived-]{derived}
\externaldocument[simplicial-]{simplicial}
\externaldocument[more-algebra-]{more-algebra}
\externaldocument[smoothing-]{smoothing}
\externaldocument[modules-]{modules}
\externaldocument[sites-modules-]{sites-modules}
\externaldocument[injectives-]{injectives}
\externaldocument[cohomology-]{cohomology}
\externaldocument[sites-cohomology-]{sites-cohomology}
\externaldocument[dga-]{dga}
\externaldocument[dpa-]{dpa}
\externaldocument[sdga-]{sdga}
\externaldocument[hypercovering-]{hypercovering}
\externaldocument[schemes-]{schemes}
\externaldocument[constructions-]{constructions}
\externaldocument[properties-]{properties}
\externaldocument[morphisms-]{morphisms}
\externaldocument[coherent-]{coherent}
\externaldocument[divisors-]{divisors}
\externaldocument[limits-]{limits}
\externaldocument[varieties-]{varieties}
\externaldocument[topologies-]{topologies}
\externaldocument[descent-]{descent}
\externaldocument[perfect-]{perfect}
\externaldocument[more-morphisms-]{more-morphisms}
\externaldocument[flat-]{flat}
\externaldocument[groupoids-]{groupoids}
\externaldocument[more-groupoids-]{more-groupoids}
\externaldocument[etale-]{etale}
\externaldocument[chow-]{chow}
\externaldocument[intersection-]{intersection}
\externaldocument[pic-]{pic}
\externaldocument[weil-]{weil}
\externaldocument[adequate-]{adequate}
\externaldocument[dualizing-]{dualizing}
\externaldocument[duality-]{duality}
\externaldocument[discriminant-]{discriminant}
\externaldocument[derham-]{derham}
\externaldocument[local-cohomology-]{local-cohomology}
\externaldocument[algebraization-]{algebraization}
\externaldocument[curves-]{curves}
\externaldocument[resolve-]{resolve}
\externaldocument[models-]{models}
\externaldocument[functors-]{functors}
\externaldocument[equiv-]{equiv}
\externaldocument[pione-]{pione}
\externaldocument[etale-cohomology-]{etale-cohomology}
\externaldocument[proetale-]{proetale}
\externaldocument[relative-cycles-]{relative-cycles}
\externaldocument[more-etale-]{more-etale}
\externaldocument[trace-]{trace}
\externaldocument[crystalline-]{crystalline}
\externaldocument[spaces-]{spaces}
\externaldocument[spaces-properties-]{spaces-properties}
\externaldocument[spaces-morphisms-]{spaces-morphisms}
\externaldocument[decent-spaces-]{decent-spaces}
\externaldocument[spaces-cohomology-]{spaces-cohomology}
\externaldocument[spaces-limits-]{spaces-limits}
\externaldocument[spaces-divisors-]{spaces-divisors}
\externaldocument[spaces-over-fields-]{spaces-over-fields}
\externaldocument[spaces-topologies-]{spaces-topologies}
\externaldocument[spaces-descent-]{spaces-descent}
\externaldocument[spaces-perfect-]{spaces-perfect}
\externaldocument[spaces-more-morphisms-]{spaces-more-morphisms}
\externaldocument[spaces-flat-]{spaces-flat}
\externaldocument[spaces-groupoids-]{spaces-groupoids}
\externaldocument[spaces-more-groupoids-]{spaces-more-groupoids}
\externaldocument[bootstrap-]{bootstrap}
\externaldocument[spaces-pushouts-]{spaces-pushouts}
\externaldocument[spaces-chow-]{spaces-chow}
\externaldocument[groupoids-quotients-]{groupoids-quotients}
\externaldocument[spaces-more-cohomology-]{spaces-more-cohomology}
\externaldocument[spaces-simplicial-]{spaces-simplicial}
\externaldocument[spaces-duality-]{spaces-duality}
\externaldocument[formal-spaces-]{formal-spaces}
\externaldocument[restricted-]{restricted}
\externaldocument[spaces-resolve-]{spaces-resolve}
\externaldocument[formal-defos-]{formal-defos}
\externaldocument[defos-]{defos}
\externaldocument[cotangent-]{cotangent}
\externaldocument[examples-defos-]{examples-defos}
\externaldocument[algebraic-]{algebraic}
\externaldocument[examples-stacks-]{examples-stacks}
\externaldocument[stacks-sheaves-]{stacks-sheaves}
\externaldocument[criteria-]{criteria}
\externaldocument[artin-]{artin}
\externaldocument[quot-]{quot}
\externaldocument[stacks-properties-]{stacks-properties}
\externaldocument[stacks-morphisms-]{stacks-morphisms}
\externaldocument[stacks-limits-]{stacks-limits}
\externaldocument[stacks-cohomology-]{stacks-cohomology}
\externaldocument[stacks-perfect-]{stacks-perfect}
\externaldocument[stacks-introduction-]{stacks-introduction}
\externaldocument[stacks-more-morphisms-]{stacks-more-morphisms}
\externaldocument[stacks-geometry-]{stacks-geometry}
\externaldocument[moduli-]{moduli}
\externaldocument[moduli-curves-]{moduli-curves}
\externaldocument[examples-]{examples}
\externaldocument[exercises-]{exercises}
\externaldocument[guide-]{guide}
\externaldocument[desirables-]{desirables}
\externaldocument[coding-]{coding}
\externaldocument[obsolete-]{obsolete}
\externaldocument[fdl-]{fdl}
\externaldocument[index-]{index}

% Theorem environments.
%
\theoremstyle{plain}
\newtheorem{theorem}[subsection]{Theorem}
\newtheorem{proposition}[subsection]{Proposition}
\newtheorem{lemma}[subsection]{Lemma}

\theoremstyle{definition}
\newtheorem{definition}[subsection]{Definition}
\newtheorem{example}[subsection]{Example}
\newtheorem{exercise}[subsection]{Exercise}
\newtheorem{situation}[subsection]{Situation}

\theoremstyle{remark}
\newtheorem{remark}[subsection]{Remark}
\newtheorem{remarks}[subsection]{Remarks}

\numberwithin{equation}{subsection}

% Macros
%
\def\lim{\mathop{\mathrm{lim}}\nolimits}
\def\colim{\mathop{\mathrm{colim}}\nolimits}
\def\Spec{\mathop{\mathrm{Spec}}}
\def\Hom{\mathop{\mathrm{Hom}}\nolimits}
\def\Ext{\mathop{\mathrm{Ext}}\nolimits}
\def\SheafHom{\mathop{\mathcal{H}\!\mathit{om}}\nolimits}
\def\SheafExt{\mathop{\mathcal{E}\!\mathit{xt}}\nolimits}
\def\Sch{\mathit{Sch}}
\def\Mor{\mathop{\mathrm{Mor}}\nolimits}
\def\Ob{\mathop{\mathrm{Ob}}\nolimits}
\def\Sh{\mathop{\mathit{Sh}}\nolimits}
\def\NL{\mathop{N\!L}\nolimits}
\def\CH{\mathop{\mathrm{CH}}\nolimits}
\def\proetale{{pro\text{-}\acute{e}tale}}
\def\etale{{\acute{e}tale}}
\def\QCoh{\mathit{QCoh}}
\def\Ker{\mathop{\mathrm{Ker}}}
\def\Im{\mathop{\mathrm{Im}}}
\def\Coker{\mathop{\mathrm{Coker}}}
\def\Coim{\mathop{\mathrm{Coim}}}

% Boxtimes
%
\DeclareMathSymbol{\boxtimes}{\mathbin}{AMSa}{"02}

%
% Macros for moduli stacks/spaces
%
\def\QCohstack{\mathcal{QC}\!\mathit{oh}}
\def\Cohstack{\mathcal{C}\!\mathit{oh}}
\def\Spacesstack{\mathcal{S}\!\mathit{paces}}
\def\Quotfunctor{\mathrm{Quot}}
\def\Hilbfunctor{\mathrm{Hilb}}
\def\Curvesstack{\mathcal{C}\!\mathit{urves}}
\def\Polarizedstack{\mathcal{P}\!\mathit{olarized}}
\def\Complexesstack{\mathcal{C}\!\mathit{omplexes}}
% \Pic is the operator that assigns to X its picard group, usage \Pic(X)
% \Picardstack_{X/B} denotes the Picard stack of X over B
% \Picardfunctor_{X/B} denotes the Picard functor of X over B
\def\Pic{\mathop{\mathrm{Pic}}\nolimits}
\def\Picardstack{\mathcal{P}\!\mathit{ic}}
\def\Picardfunctor{\mathrm{Pic}}
\def\Deformationcategory{\mathcal{D}\!\mathit{ef}}
\title{A trivial Kan fibration has the lifting property against termwise injective maps}
\author{The Stacks Project Authors}
\date{stacks-project:simplicial.tex:local-snapshot-20261012; GNU FDL 1.2 or later}
\begin{document}
\maketitle
\noindent\textit{Editorial note.} The counted statement is the author's own lemma environment labelled lemma-trivial-kan (source lines 5767-5779); the proof environment begins on the second line after it and is its own proof, closing at source line 5800.\par\medskip
\noindent\textit{Editorial note.} The statement and the proof invoke two author results of the same chapter by label; both are retained as context spans and the generator expands the context transitively (budget 800 lines for this card), while reference targets in other chapters that carry canonical project tags are auto-declared as non-load-bearing source pointers.\par\medskip
\noindent\textit{Editorial note.} The retained preamble is the chapter's own minimal preamble (sources/stacks-simplicial/preamble.tex, taken with its real line count), a verbatim excerpt of the native Xy-pic setup, the native theorem declarations with the mathematical macros, and the project's own no-op reference environments; no macro was invented or redefined.\par\medskip
\noindent\textit{Editorial scope.} One original result of the Stacks Project's Simplicial Methods chapter, counted once: a trivial Kan fibration lifts every solid diagram whose horizontal map is termwise injective.\par\medskip
\noindent\textit{Allowed standard prerequisites (editorial).}\begin{itemize}
\item Simplicial sets and termwise injective maps. The category of simplicial sets, the standard simplices Delta[n] and their boundaries, and termwise injectivity of a morphism of simplicial sets.. It is the setting and the hypothesis of the statement (source lines 5769-5777).
\item The Yoneda description of simplices. A simplex x of a simplicial set W corresponds to a morphism Delta[n] to W (the author's lemma-simplex-map, retained as context).. It produces the map varphi : Delta[n] to Z' attached to the chosen x (source lines 5785-5787).
\item The pushout describing an elementary extension. The simplicial subset generated by Z, a nondegenerate simplex x and its degeneracies is the pushout of Z and Delta[n] along the boundary (the author's lemma-glue-simplex, retained as context).. It lets b and beta be glued into a morphism on Z' (source lines 5792-5794).
\item The Kan lifting condition and Zorn's lemma. Trivial Kan fibrations lift horns against boundary inclusions, and Zorn's lemma supplies a maximal extension.. The lifting condition gives beta, and Zorn finishes the proof (source lines 5790-5791, 5799).
\end{itemize}\medskip
\begin{lemma}
\label{lemma-simplex-map}
Let $U$ be a simplicial set. Let $n \geq 0$ be an integer.
There is a canonical bijection
$$
\Mor(\Delta[n], U)
\longrightarrow
U_n
$$
which maps a morphism $\varphi$ to the value of $\varphi$
on the unique nondegenerate $n$-simplex of $\Delta[n]$.
\end{lemma}

\begin{proof}
Omitted.
\end{proof}

\begin{lemma}
\label{lemma-glue-simplex}
Let $U \subset V$ be simplicial sets.
Suppose $n \geq 0$ and $x \in V_n$, $x \not \in U_n$ are such that
\begin{enumerate}
\item $V_i = U_i$ for $i < n$,
\item $V_n = U_n \cup \{x\}$,
\item any $z \in V_j$, $z \not \in U_j$ for $j > n$
is degenerate.
\end{enumerate}
Let $\Delta[n] \to V$ be the unique morphism mapping the
nondegenerate $n$-simplex of $\Delta[n]$ to $x$.
In this case the diagram
$$
\xymatrix{
\Delta[n] \ar[r] & V \\
i_{(n - 1)!} \text{sk}_{n - 1} \Delta[n] \ar[r] \ar[u] & U \ar[u]
}
$$
is a pushout diagram.
\end{lemma}

\begin{proof}
Let us denote $\partial \Delta[n] = i_{(n - 1)!} \text{sk}_{n - 1} \Delta[n]$
for convenience. There is a natural map
$U \amalg_{\partial \Delta[n]} \Delta[n] \to V$.
We have to show that it is bijective in degree $j$
for all $j$. This is clear for $j \leq n$. Let $j > n$.
The third condition means that any $z \in V_j$, $z \not \in U_j$
is a degenerate simplex, say $z = s^{j - 1}_i(z')$. Of course
$z' \not \in U_{j - 1}$. By induction it follows that $z'$
is a degeneracy of $x$. Thus we conclude that all $j$-simplices
of $V$ are either in $U$ or degeneracies of $x$. This implies
that the map $U \amalg_{\partial \Delta[n]} \Delta[n] \to V$
is surjective. Note that a nondegenerate simplex of
$U \amalg_{\partial \Delta[n]} \Delta[n]$ is either
the image of a nondegenerate simplex of $U$, or
the image of the (unique) nondegenerate $n$-simplex
of $\Delta[n]$. Since clearly $x$ is nondegenerate we
deduce that $U \amalg_{\partial \Delta[n]} \Delta[n] \to V$
maps nondegenerate simplices to nondegenerate simplices
and is injective on nondegenerate simplices. Hence it is
injective, by Lemma \ref{lemma-injective-map-simplicial-sets}.
\end{proof}

\begin{lemma}
\label{lemma-injective-map-simplicial-sets}
Let $f : U \to V$ be a morphism of simplicial sets.
Suppose that (a) the image of every nondegenerate simplex of
$U$ is a nondegenerate simplex of $V$ and (b) the restriction
of $f$ to a map from the set of nondegenerate simplices of $U$
to the set of nondegenerate simplices of $V$ is injective.
Then $f_n$ is injective for all $n$.
Same holds with ``injective'' replaced by
``surjective'' or ``bijective''.
\end{lemma}

\begin{proof}
Under hypothesis (a) we see that the map $f$ preserves
the disjoint union decompositions of the splitting
of Lemma \ref{lemma-splitting-simplicial-sets}, in other words
that we get commutative diagrams
$$
\xymatrix{
\coprod\nolimits_{\varphi : [n] \to [m]\text{ surjective}}
N(U_m)
\ar[r] \ar[d] &
U_n \ar[d] \\
\coprod\nolimits_{\varphi : [n] \to [m]\text{ surjective}}
N(V_m)
\ar[r] &
V_n.
}
$$
And then (b) clearly shows that the left vertical arrow is
injective (resp.\ surjective, resp.\ bijective).
\end{proof}

\begin{lemma}
\label{lemma-splitting-simplicial-sets}
Let $U$ be a simplicial set. Then $U$ has a unique splitting
with $N(U_m)$ equal to the set of nondegenerate $m$-simplices.
\end{lemma}

\begin{proof}
From the definition it follows immediately, that if there is a
splitting then $N(U_m)$ has to be the set of nondegenerate simplices.
Let $x \in U_n$. Suppose that there are surjections $\varphi : [n] \to [k]$
and $\psi : [n] \to [l]$ and nondegenerate simplices
$y \in U_k$, $z \in U_l$ such that $x = U(\varphi)(y)$
and $x = U(\psi)(z)$. Choose a right inverse $\xi : [l] \to [n]$
of $\psi$, i.e., $\psi \circ \xi = \text{id}_{[l]}$.
Then $z = U(\xi)(x)$. Hence $z = U(\xi)(x) = U(\varphi \circ \xi)(y)$.
Since $z$ is nondegenerate we conclude that $\varphi \circ \xi :
[l] \to [k]$ is surjective, and hence $l \geq k$. Similarly
$k \geq l$. Hence we see that $\varphi \circ \xi : [l] \to [k]$
has to be the identity map for any choice of right inverse
$\xi$ of $\psi$. This easily implies that $\psi = \varphi$.
\end{proof}

\begin{lemma}
\label{lemma-trivial-kan}
Let $f : X \to Y$ be a trivial Kan fibration of simplicial sets.
For any solid commutative diagram
$$
\xymatrix{
Z \ar[r]_b \ar[d] & X \ar[d] \\
W \ar[r]^a \ar@{-->}[ru] & Y
}
$$
of simplicial sets with $Z \to W$ (termwise) injective
a dotted arrow exists making the diagram commute.
\end{lemma}

\begin{proof}
Suppose that $Z \not = W$. Let $n$ be the smallest integer such that
$Z_n \not = W_n$. Let $x \in W_n$, $x \not \in Z_n$. Denote
$Z' \subset W$ the simplicial subset containing $Z$, $x$, and
all degeneracies of $x$. Let
$\varphi : \Delta[n] \to Z'$ be the morphism corresponding to $x$
(Lemma \ref{lemma-simplex-map}).
Then $\varphi|_{\partial \Delta[n]}$ maps into $Z$
as all the nondegenerate simplices of $\partial \Delta[n]$ end up in $Z$.
By assumption we can extend $b \circ \varphi|_{\partial \Delta[n]}$
to $\beta : \Delta[n] \to X$.
By Lemma \ref{lemma-glue-simplex} the simplicial set $Z'$
is the pushout of $\Delta[n]$ and $Z$ along $\partial \Delta[n]$.
Hence $b$ and $\beta$ define a morphism $b' : Z' \to X$.
In other words, we have extended the morphism $b$ to a bigger
simplicial subset of $Z$.

\medskip\noindent
The proof is finished by an application of Zorn's lemma (omitted).
\end{proof}

\end{document}
